\documentclass[fleqn,landscape]{article}

%% Created with wxMaxima 20.06.6

\setlength{\parskip}{\medskipamount}
\setlength{\parindent}{0pt}
\usepackage{iftex}
\ifPDFTeX
  % PDFLaTeX or LaTeX 
  \usepackage[utf8]{inputenc}
  \usepackage[T1]{fontenc}
  \DeclareUnicodeCharacter{00B5}{\ensuremath{\mu}}
\else
  %  XeLaTeX or LuaLaTeX
  \usepackage{fontspec}
\fi
\usepackage{graphicx}
\usepackage{color}
\usepackage{amsmath}
\usepackage{grffile}
\usepackage{ifthen}
\newsavebox{\picturebox}
\newlength{\pictureboxwidth}
\newlength{\pictureboxheight}
\newcommand{\includeimage}[1]{
    \savebox{\picturebox}{\includegraphics{#1}}
    \settoheight{\pictureboxheight}{\usebox{\picturebox}}
    \settowidth{\pictureboxwidth}{\usebox{\picturebox}}
    \ifthenelse{\lengthtest{\pictureboxwidth > .95\linewidth}}
    {
        \includegraphics[width=.95\linewidth,height=.80\textheight,keepaspectratio]{#1}
    }
    {
        \ifthenelse{\lengthtest{\pictureboxheight>.80\textheight}}
        {
            \includegraphics[width=.95\linewidth,height=.80\textheight,keepaspectratio]{#1}
            
        }
        {
            \includegraphics{#1}
        }
    }
}
\newlength{\thislabelwidth}
\DeclareMathOperator{\abs}{abs}
\usepackage{animate} % This package is required because the wxMaxima configuration option
                      % "Export animations to TeX" was enabled when this file was generated.

\definecolor{labelcolor}{RGB}{100,0,0}

\begin{document}

\pagebreak{}
{\Huge {\scshape Order of operations in SWIRL}}
\setcounter{section}{0}
\setcounter{subsection}{0}
\setcounter{figure}{0}

\textbf{Affine Transformations} are \textit{composable}. \textbf{SWIRL} performs them in the following order:
\begin{itemize}
	\item{P} pivoting
	\item{K} skewing
	\item{S} scaling
	\item{R} rotation
	\item{T} translation
\end{itemize}
The mnemonic ``\emph{PiKe SoRTie}'' simplifies memorizing it. Here we show the forms of these elemental transforms. When multiplied by a bi-dimensional \textit{row vector} in the form $\left( x \; y \; 1 \right)$, the product will yield the indicated result. Furthermore, the matrices can be multiplied together, so the vector only needs to be multiplied against the final matrix product.


\noindent
%%%%%%%%
%% INPUT:
\begin{minipage}[t]{4.000000em}\color{red}\bfseries
(\% i1) 
\end{minipage}
\begin{minipage}[t]{\textwidth}\color{blue}
P: matrix([1, 0, 0], [0, 1, 0], [-p\_x, -p\_y, 1]);
\end{minipage}
%%%% OUTPUT:
\[\displaystyle \tag{\% o1} 
\begin{pmatrix}1 & 0 & 0\\
0 & 1 & 0\\
-{p_x} & -{p_y} & 1\end{pmatrix}\mbox{}
\]
%%%%%%%%%%%%%%%%


\noindent
%%%%%%%%
%% INPUT:
\begin{minipage}[t]{4.000000em}\color{red}\bfseries
(\% i2)
\end{minipage}
\begin{minipage}[t]{\textwidth}\color{blue}
K: matrix([1, k\_y, 0], [k\_x, 1, 0], [0, 0, 1]);
\end{minipage}
%%%% OUTPUT:
\[\displaystyle \tag{\% o2} 
\begin{pmatrix}1 & {k_y} & 0\\
{k_x} & 1 & 0\\
0 & 0 & 1\end{pmatrix}\mbox{}
\]
%%%%%%%%%%%%%%%%


\noindent
%%%%%%%%
%% INPUT:
\begin{minipage}[t]{4.000000em}\color{red}\bfseries
(\% i3)
\end{minipage}
\begin{minipage}[t]{\textwidth}\color{blue}
S: matrix([s\_x, 0, 0], [0, s\_y, 0], [0,0,1]);
\end{minipage}
%%%% OUTPUT:
\[\displaystyle \tag{\% o3} 
\begin{pmatrix}{s_x} & 0 & 0\\
0 & {s_y} & 0\\
0 & 0 & 1\end{pmatrix}\mbox{}
\]
%%%%%%%%%%%%%%%%


\noindent
%%%%%%%%
%% INPUT:
\begin{minipage}[t]{4.000000em}\color{red}\bfseries
(\% i4)
\end{minipage}
\begin{minipage}[t]{\textwidth}\color{blue}
R: matrix([cos(\ensuremath{\theta}), sin(\ensuremath{\theta}), 0], [-sin(\ensuremath{\theta}), cos(\ensuremath{\theta}), 0], [0, 0, 1]);
\end{minipage}
%%%% OUTPUT:
\[\displaystyle \tag{\% o4} 
\begin{pmatrix}\cos{\left( \theta \right) } & \sin{\left( \theta \right) } & 0\\
-\sin{\left( \theta \right) } & \cos{\left( \theta \right) } & 0\\
0 & 0 & 1\end{pmatrix}\mbox{}
\]
%%%%%%%%%%%%%%%%


\noindent
%%%%%%%%
%% INPUT:
\begin{minipage}[t]{4.000000em}\color{red}\bfseries
(\% i5)
\end{minipage}
\begin{minipage}[t]{\textwidth}\color{blue}
T: matrix([1, 0, 0], [0, 1, 0], [t\_x, t\_y, 1]);
\end{minipage}
%%%% OUTPUT:
\[\displaystyle \tag{\% o5} 
\begin{pmatrix}1 & 0 & 0\\
0 & 1 & 0\\
{t_x} & {t_y} & 1\end{pmatrix}\mbox{}
\]
%%%%%%%%%%%%%%%%
Now for the cumulative transformations. Note that the "." operator means matrix multiplication in Maxima:


\noindent
%%%%%%%%
%% INPUT:
\begin{minipage}[t]{4.000000em}\color{red}\bfseries
(\% i6)
\end{minipage}
\begin{minipage}[t]{\textwidth}\color{blue}
P.K;
\end{minipage}
%%%% OUTPUT:
\[\displaystyle \tag{\% o6} 
\begin{pmatrix}1 & {k_y} & 0\\
{k_x} & 1 & 0\\
-{k_x} {p_y}-{p_x} & -{p_y}-{k_y} {p_x} & 1\end{pmatrix}\mbox{}
\]
%%%%%%%%%%%%%%%%


\noindent
%%%%%%%%
%% INPUT:
\begin{minipage}[t]{4.000000em}\color{red}\bfseries
(\% i7)
\end{minipage}
\begin{minipage}[t]{\textwidth}\color{blue}
P.K.S;
\end{minipage}
%%%% OUTPUT:
\[\displaystyle \tag{\% o7} 
\begin{pmatrix}{s_x} & {k_y} {s_y} & 0\\
{k_x} {s_x} & {s_y} & 0\\
-{k_x} {p_y} {s_x}-{p_x} {s_x} & -{p_y} {s_y}-{k_y} {p_x} {s_y} & 1\end{pmatrix}\mbox{}
\]
%%%%%%%%%%%%%%%%


\noindent
%%%%%%%%
%% INPUT:
\begin{minipage}[t]{4.000000em}\color{red}\bfseries
(\% i8)
\end{minipage}
\begin{minipage}[t]{\textwidth}\color{blue}
P.K.S.R;
\end{minipage}
%%%% OUTPUT:
\[\displaystyle \tag{\% o8}
\small{
\begin{pmatrix}{s_x} \cos{\left( \theta \right) }-{k_y} {s_y} \sin{\left( \theta \right) } & {s_x} \sin{\left( \theta \right) }+{k_y} {s_y} \cos{\left( \theta \right) } & 0\\
{k_x} {s_x} \cos{\left( \theta \right) }-{s_y} \sin{\left( \theta \right) } & {k_x} {s_x} \sin{\left( \theta \right) }+{s_y} \cos{\left( \theta \right) } & 0\\
-{p_x} \left( {s_x} \cos{\left( \theta \right) }-{k_y} {s_y} \sin{\left( \theta \right) }\right) -{p_y} \left( {k_x} {s_x} \cos{\left( \theta \right) }-{s_y} \sin{\left( \theta \right) }\right)  & -{p_y} \left( {k_x} {s_x} \sin{\left( \theta \right) }+{s_y} \cos{\left( \theta \right) }\right) -{p_x} \left( {s_x} \sin{\left( \theta \right) }+{k_y} {s_y} \cos{\left( \theta \right) }\right)  & 1\end{pmatrix}\mbox{}
}
\]
%%%%%%%%%%%%%%%%


\noindent
%%%%%%%%
%% INPUT:
\begin{minipage}[t]{4.000000em}\color{red}\bfseries
(\% i9)
\end{minipage}
\begin{minipage}[t]{\textwidth}\color{blue}
P.K.S.R.T;
\end{minipage}
%%%% OUTPUT:
\[\displaystyle \tag{\% o9}
\small{ 
\begin{pmatrix}{s_x} \cos{\left( \theta \right) }-{k_y} {s_y} \sin{\left( \theta \right) } & {s_x} \sin{\left( \theta \right) }+{k_y} {s_y} \cos{\left( \theta \right) } & 0\\
{k_x} {s_x} \cos{\left( \theta \right) }-{s_y} \sin{\left( \theta \right) } & {k_x} {s_x} \sin{\left( \theta \right) }+{s_y} \cos{\left( \theta \right) } & 0\\
-{p_x} \left( {s_x} \cos{\left( \theta \right) }-{k_y} {s_y} \sin{\left( \theta \right) }\right) -{p_y} \left( {k_x} {s_x} \cos{\left( \theta \right) }-{s_y} \sin{\left( \theta \right) }\right) +{t_x} & -{p_y} \left( {k_x} {s_x} \sin{\left( \theta \right) }+{s_y} \cos{\left( \theta \right) }\right) -{p_x} \left( {s_x} \sin{\left( \theta \right) }+{k_y} {s_y} \cos{\left( \theta \right) }\right) +{t_y} & 1\end{pmatrix}\mbox{}
}
\]
%%%%%%%%%%%%%%%%
This yields the complete Local Transformation. The [final] World Transformation is obtained when the LocalTransformation is multiplied by its container's World Transformation. §
\end{document}
